\clearpage
\section{Validation}

Validation is based on automated \textbf{JUnit 5} tests.
The goal is not only to check that one implementation returns plausible numbers, but also to verify that the three paradigms preserve the same \textbf{report semantics}.

The test suite is organized as follows:

\begin{itemize}[leftmargin=*]
  \item \textbf{\texttt{FSStatTest}}: unit coverage for the shared model, parsing, and the three engines on a small fixture;
  \item \textbf{\texttt{FSStatRobustnessCorrectnessTest}}: edge cases and cross-implementation validation on identical datasets;
  \item \textbf{\texttt{FSStatCLITest}}: CLI parsing, dispatching, and reactive completion checks.
\end{itemize}

\subsection{Correctness Coverage}

The suite checks the main functional cases required by the assignment:

\begin{itemize}[leftmargin=*]
  \item \textbf{empty directories}: all counters must be zero;
  \item \textbf{flat directories}: known file sizes must be assigned to the expected bands;
  \item \textbf{nested directories}: recursive traversal must not lose files;
  \item \textbf{wide/deep deterministic trees}: larger fixtures must produce stable totals and histograms;
  \item \textbf{boundary values}: files at \texttt{0}, band limits,
  \texttt{MaxFS}, and above \texttt{MaxFS} must be classified deterministically;
  \item \textbf{large files}: sparse files larger than 32-bit integer ranges must still be classified correctly.
\end{itemize}

\subsection{Cross-Implementation Oracle}

The main correctness strategy is \textbf{comparative}.
For deterministic fixtures, the Virtual Threads result is used as a reference report and is compared with the Reactive and Event-Loop results.
The comparison checks:

\begin{itemize}[leftmargin=*]
  \item the \textbf{total file count};
  \item the length of the \textbf{histogram};
  \item every individual \textbf{band counter}.
\end{itemize}

This is useful because the implementations have different control structures.
If the imperative virtual-thread version, the RxJava stream version, and the Vert.x callback version agree on the same non-trivial tree, the risk of a paradigm-specific scheduling error is lower.

The tests also validate the shared model directly: band boundaries, duration formatting, size-unit parsing, and formatted labels are checked independently of the filesystem traversal.

\subsection{Robustness Coverage}

The robustness tests cover filesystem situations that may occur in practice:

\begin{itemize}[leftmargin=*]
  \item \textbf{invalid roots}: non-existing paths must report an error;
  \item \textbf{wrong root type}: a regular file cannot be used as the scan root;
  \item \textbf{permission issues}: unreadable entries must not crash the scan, where the host filesystem enforces permissions;
  \item \textbf{symbolic-link cycles}: repeated logical directories must not cause infinite traversal, where symbolic links are available;
  \item \textbf{concurrent mutation}: files deleted during a scan must not make the program crash;
  \item \textbf{special entries}: non-regular nodes are skipped, where the platform exposes them.
\end{itemize}

\subsection{Cancellation and Interface Checks}

The interactive extension is validated through cancellation and interface tests:

\begin{itemize}[leftmargin=*]
  \item the Event-Loop cancellation test cancels immediately after the
  \texttt{FSReportJob} is returned and checks that no terminal callback is emitted afterward;
  \item the CLI tests verify argument parsing, size-unit aliases, paradigm selection, and dispatching to \texttt{vt}, \texttt{rx}, and \texttt{loop};
  \item the reactive CLI path checks that the final listener completion uses the latest report emitted by the \texttt{Observable}.
\end{itemize}

The Event-Loop cancellation test avoids arbitrary sleeps and machine-speed assumptions.
This makes it more stable on local machines and CI runners.

\subsection{Platform-Dependent Checks}

Some checks depend on the host environment:

\begin{itemize}[leftmargin=*]
  \item permission changes may not be applied by the local filesystem;
  \item symbolic links may be unavailable or restricted;
  \item special files such as \texttt{/dev/null} are Unix-specific.
\end{itemize}

When one of these features is not available, the corresponding test skips only the platform-specific assertion.
The core consistency checks remain valid.

The final verification run executed \texttt{26} tests and all of them passed.
